\section{Los dos grandes paradigmas algorítmicos de Deep RL}

Después de presentar los conceptos básicos de RL, esta sección se centra en cómo implementar algoritmos de aprendizaje por refuerzo usando \textbf{redes neuronales profundas}. Deep RL tiene dos grandes paradigmas:

\begin{figure}[H]
\centering
\includegraphics[width=0.85\textwidth]{slides-images/slide-20.jpg}
\caption{Los dos grandes paradigmas algorítmicos de Deep RL: Value Learning y Policy Learning\protect\footnotemark}
\end{figure}
\footnotetext{Diapositiva 20 de las láminas.}

\begin{importantbox}{Value Learning vs. Policy Learning}
\begin{itemize}
    \item \textbf{Value Learning}(aprendizaje de valor): primero se aprende la Q-function $Q(s, a)$, y luego se determina la acción óptima mediante $a = \arg\max_a Q(s,a)$
    \item \textbf{Policy Learning}(aprendizaje de política): se aprende directamente la función de política $\pi(s)$, y se muestrea la acción a partir de la política $a \sim \pi(s)$
\end{itemize}
Los dos métodos tienen ventajas y desventajas distintas, y son adecuados para diferentes tipos de problemas.
\end{importantbox}

\subsection{Resumen de la sección}
La idea central de Deep RL es usar redes neuronales profundas como aproximador de funciones para aprender la Q-function o la política. Value Learning resuelve el problema de decisión de manera indirecta (primero se aprende $Q$ y luego se decide), mientras que Policy Learning lo resuelve de manera directa (se aprende la política directamente). Las siguientes dos secciones presentarán en detalle, respectivamente, los algoritmos representativos de estos dos paradigmas.
